\chapter{Scale Analysis and Asymptotic Reduction}\label{ch:4}

\begin{chaptersummary}
Kinetic, continuum, and electromagnetic models contain several scales and mechanisms. Selecting among them begins with nondimensionalization, which expresses relative magnitudes as scale ratios. Asymptotic analysis then constructs limiting models and boundary layers, locating the validity of the preceding model families and guiding numerical resolution and error analysis.
\end{chaptersummary}
\learninggoals{Construct dimensionless parameters from characteristic scales; determine inner and outer regions by dominant balance; identify the fast scales removed by quasineutral, strong-field, and low-frequency reductions.}
\studyroute{The construction of matching conditions and proofs of uniform error estimates are omitted.}{singular perturbation; dominant balance; matched asymptotic expansions}{For $\epsilon u''+u'=0$ with $u(0)=0,u(1)=1$, construct inner and outer solutions, identify the boundary-layer location, and obtain a uniform approximation.}

\section{Nondimensionalization, Scale Ordering, and Asymptotic Expansions}\index{Singular perturbation}
Common dimensionless groups include
\[
\mathrm{Kn}=\frac{\lambda}{L},
\qquad
\mathrm{Re}=\frac{\rho UL}{\mu},
\qquad
\mathrm{Ma}=\frac{U}{c}.
\]
These parameters compare the magnitudes of different physical mechanisms. Model selection determines which terms remain in a specified limit.

\begin{authornote}
\textbf{Potential Question:}
"Why discuss dimensionless groups immediately after continuum equations? How do these numbers determine model selection?"

\textbf{Author Supplement:}
Numerical magnitudes of dimensional quantities depend on the chosen units and cannot be compared directly. After characteristic length, velocity, and time scales are selected, nondimensionalization expresses relative term sizes through dimensionless parameters. Their magnitudes indicate potentially dominant mechanisms and limits admitting reduced models. The scales themselves must be chosen for the particular problem.

\textbf{Core Conceptual Connection:}
\[
\text{continuum equations}
\longrightarrow
\text{dimensionless ratios}
\longrightarrow
\text{model regimes}
\]
\end{authornote}

Taylor expansion describes local variation through derivatives; asymptotic expansion orders scales in a small-parameter limit. Begin with the Taylor form:
\[
f(x+h)
=
f(x)+hf'(x)+\frac{h^2}{2}f''(x)+\cdots.
\]

\begin{askai}
Ask AI:
\begin{quote}
Compare Taylor and asymptotic expansions. Give an example of a useful divergent asymptotic series and explain why a smaller remainder at each fixed truncation order can be more relevant than convergence of the infinite sum.
\end{quote}

Follow-up:
\begin{quote}
For a simple function, calculate truncation errors at different orders and illustrate optimal truncation near the smallest term.
\end{quote}
\end{askai}

Taylor expansion is particularly useful for analyzing local differences between discrete and continuous operators.

An asymptotic expansion is more general. As \(\epsilon\to0\), suppose the scale functions satisfy
\[
\phi_{n+1}(\epsilon)=o(\phi_n(\epsilon)),
\]
and, for every fixed \(N\),
\[
f(\epsilon)
-
\sum_{n=0}^{N}a_n\phi_n(\epsilon)
=
o(\phi_N(\epsilon)),
\]
Then write
\[
f(\epsilon)\sim\sum_{n=0}^{\infty}a_n\phi_n(\epsilon).
\]

At each fixed truncation order, an asymptotic expansion requires a remainder smaller than the last retained scale. The series itself may diverge, and its scales may include fractional powers, logarithms, and exponentials:
\[
1,\quad
\epsilon^{1/2},\quad
\epsilon\log\epsilon,\quad
\ee^{-1/\epsilon}.
\]

A Taylor series is thus a special expansion generated by local derivatives in ordinary powers, whereas an asymptotic expansion emphasizes an ordering of scales.

Rapidly growing coefficients can cause an asymptotic series to improve initially and then deteriorate. For fixed \(\epsilon\), the terms
\[
|a_n\epsilon^n|
\]
may first decrease and then increase. Truncate near the smallest term rather than pursuing an infinite sum; this is optimal truncation.

If substituting
\[
u=u_0+\epsilon u_1+\epsilon^2u_2+\cdots
\]
preserves the differential order and boundary-condition structure at successive orders, the perturbation is regular.

If a small parameter multiplies the highest derivative, for example,
\[
\epsilon u''+u'=0,
\]
the limit \(\epsilon\to0\) lowers the order, and the outer solution cannot satisfy all boundary conditions. This is a singular perturbation, typically requiring boundary layers, matched asymptotic expansions, multiple scales, or WKB methods.

\section{Plasma Example: Small Debye Length and the Quasineutral Limit}

Consider the nondimensional one-dimensional Poisson equation
\[
-\epsilon^2\frac{\partial^2\phi}{\partial x^2}
=
n_i-n_e,
\qquad
\epsilon=\frac{\lambda_D}{L}\ll1.
\]
Here \(\lambda_D\) is the Debye length and \(L\) the macroscopic scale.

Setting \(\epsilon=0\) directly gives
\[
n_i=n_e,
\]
the quasineutrality condition. Net charge density is a higher-order quantity in the bulk.

If the boundary potential is incompatible with the bulk quasineutral solution, the small coefficient of the highest derivative makes the problem singular. An outer expansion in the bulk is
\[
\phi^{\mathrm{out}}
=
\phi_0+\epsilon\phi_1+\epsilon^2\phi_2+\cdots,
\]
For a layer with $O(1)$ potential variation and normalized net charge, the balance $\epsilon^2/\delta^2=O(1)$ gives $\delta=O(\epsilon)$. Under this ordering, introduce a stretched coordinate near the boundary:
\[
\xi=\frac{x}{\epsilon}.
\]

\begin{authornote}
\textbf{Potential Question:}
"Why introduce \(\xi=x/\epsilon\) near the boundary rather than an arbitrary new coordinate?"

\textbf{Author Supplement:}
A small coefficient can make the highest-derivative term negligible in the bulk, but the reduced equation may not satisfy all boundary conditions. Near a boundary, rapid variation over a short distance restores that derivative term to the same order as the others. Estimate all terms using an unknown layer thickness, determine it by dominant balance, and only then introduce the stretched coordinate.

\textbf{Core Conceptual Connection:}
\[
\text{singular perturbation}
\longrightarrow
\text{boundary-layer scale}
\longrightarrow
\xi=x/\epsilon
\]
\end{authornote}

At this scale,
\[
\frac{\partial}{\partial x}
=
\frac{1}{\epsilon}\frac{\partial}{\partial\xi},
\]

\begin{askai}
Ask AI:
\begin{quote}
Introduce an unknown boundary-layer width \(\delta(\epsilon)\) and determine how \(\delta\) scales with \(\epsilon\) by dominant balance. Then compare the stretched coordinate \(x/\epsilon\) with that scale and explain \(\xi=x/\epsilon\).
\end{quote}

Follow-up:
\begin{quote}
Draw a conceptual diagram showing the outer solution, inner solution, and overlap region, and identify where quasineutrality is valid.
\end{quote}
\end{askai}

The Poisson equation becomes
\[
-\frac{\partial^2\phi^{\mathrm{in}}}{\partial\xi^2}
=
n_i-n_e.
\]

Under this ordering, the Debye-scale charge-separation layer matches the wall potential to the bulk solution. Its profile also requires particle response, fluxes, and wall conditions. Other potential amplitudes or transport orderings may introduce additional transition layers.

This example illustrates characteristic features of asymptotic analysis:
\begin{enumerate}
  \item The small parameter is an explicit scale ratio, \(\lambda_D/L\).
  \item Bulk and boundary-layer regions use different coordinates.
  \item Setting the small parameter to zero lowers the differential order.
  \item Inner and outer solutions must match in an overlap region.
  \item The limiting model depends on the spatial region and boundary conditions.
\end{enumerate}

Related structures occur in low-frequency and strong-field plasma limits, drift kinetics, gyrokinetics, and MHD. Identify the small parameter, variables ordered at the same magnitude, any reduction in differential order, and any fast-time, short-space, or boundary-layer scales.

Common asymptotic notation has the following meanings:
\[
f=g+\Order(\epsilon^p)
\]
means that the remainder is bounded by a constant times \(\epsilon^p\);
\[
f=g+o(\epsilon^p)
\]
means that the remainder divided by \(\epsilon^p\) tends to zero;
\[
f\sim g
\]
may mean \(f/g\to1\) or an order-by-order asymptotic expansion, depending on context.

\section{Strong Fields, Low Frequencies, and Gyroscales}
A charged particle's gyrofrequency and characteristic gyroradius are
\[
\Omega_s=\frac{|q_s|B_0}{m_s},\qquad
\rho_{Ls}=\frac{v_{\mathrm{th},s}}{\Omega_s}.
\]
When $\rho_{Ls}/L\ll1$ and $\omega/\Omega_s\ll1$, variables can be organized into fast gyrophase and slow guiding-center motion.
Drift kinetics further assumes $k_\perp\rho_{Ls}\ll1$; gyrokinetics permits perturbations with $k_\perp\rho_{Ls}=O(1)$,
retaining finite-Larmor-radius effects. Chapter~\ref{ch:19} develops the five-dimensional distribution, gyroaveraging, and self-consistent field equations.
